\section{Baseline: RSIAT}
\label{sec:background}

RSIAT~\cite{rsiat} uses a ViT-B/16 pre-trained on ImageNet-21k~\cite{dosovitskiy2021vit}
with a shared bottleneck adapter (width 64) and a cosine classifier with scale
$s=50$. Training proceeds as follows.

\begin{itemize}
  \item \emph{Task $0$}: the adapter and classifier are trained with
    cross-entropy and a representation-steering loss.
  \item \emph{Each task $t>0$}: the adapter is trained on the new classes, with
    a feature-alignment loss to the previous model through a learned residual
    projector.
  \item \emph{Statistics}: after each task, a mean $\boldsymbol\mu_c$ and a
    covariance $\Sigma_c$ (plus $10^{-3}I$) are stored for every new class.
    They are computed on its training images with test-time transforms.
  \item \emph{Drift compensation (SSCA)}: the old means are shifted by the
    feature drift estimated on the new data. The covariances are never
    updated.
  \item \emph{Classifier alignment (CA)}: the classifier is re-trained for 10
    epochs on 256 Gaussian pseudo-features per class drawn from
    $\mathcal N(\gamma_t\boldsymbol\mu_c,\Sigma_c)$. Here $\gamma_t$ is a
    task-dependent mean shrinkage.
\end{itemize}

At inference, a test feature $\mathbf z\in\mathbb R^{768}$ is assigned to
$\arg\max_c\, \ell_c(\mathbf z)$ with $\ell_c=s\cos(\mathbf w_c,\mathbf z)$. We call
this the \emph{linear head}. All our methods change only this last step and
leave training untouched.
